% Chapter: what the available writing tools teach us
\chapter{Open-source tools}
\label{sec:oss}

One of our rejected surveys earned the best score in a small test of writing
tools. The program found few stock transitions, repeated openings, or phrases
associated with machine-written prose. It gave the survey 10 points out of
100, where a lower score meant a cleaner page. A human reader still could not
follow it.

That result is a useful place to begin. It shows both why writing tools belong
in Humanvoice and why none of them can judge the finished document. A tool can
notice a repeated phrase. It can count sentence lengths. With the right parser,
it can point to a line in a LaTeX source file. It cannot tell whether an
economist has explained the mechanism, whether an executive can reconstruct
the decision, or whether a paragraph has asked the reader to absorb six new
ideas at once.

We examined the open-source field with that boundary in mind. The detailed
package records, including versions, licences, repository activity, and the
reasons for each disposition, appear in Appendix~\ref{app:oss-records}. This
chapter develops the smaller number of ideas that change the product.

\section{A low score and a failed reader}
\label{sec:oss-calibration}

The test used sloptrim, a local program that counts a documented set of
phrasing and structural patterns. We gave it three documents whose histories
we already knew. Two were versions of the same survey on zero-lower-bound
models. The first was an internal working note. The second had been rewritten
and accepted by its owner. The third was a technically serious survey that a
reader had rejected because it did not teach.

\begin{table}[H]
\centering
\small
\begin{tabular}{@{}lccl@{}}
\toprule
Document & Score & Band & Human verdict \\
\midrule
ZLB survey, internal markdown & 45 & mixed & rejected (register) \\
ZLB survey, final LaTeX & 25 & light tells & \textbf{accepted} \\
Dense technical survey & 10 & clean & rejected (does not teach) \\
\bottomrule
\end{tabular}
\caption{sloptrim v0.9.2 scores for three documents with recorded reader
outcomes. Lower scores indicate fewer of the patterns the program measures.}
\label{tab:calibration}
\end{table}

The first comparison went as hoped. The score fell from 45 to 25 after the ZLB
survey lost its internal labels, crowded inline headings, and conspicuously
regular phrasing. A one-second check recovered a change that had taken a human
several readings to describe. Repeating that check after a later edit could
warn us that the old habits had returned.

The third row changes the interpretation. The dense survey looked cleaner to
the program than the accepted document. It did not contain many of the phrases
or rhythms the program had been written to count. Its failure lay elsewhere.
Each paragraph introduced too many distinctions, and the relations among them
remained implicit. The reader had to unpack the argument while learning it.

The program did not make an error. It answered a narrow question accurately:
how often does this text exhibit these observable patterns? The reader answered
a different question: can I understand the argument well enough to use it?
Humanvoice needs both observations, but they must never be confused.

\hvFigureScoreBoundary

This calibration also explains why a single readability number would be a
poor product. A score is useful when it leads back to evidence a writer can
inspect. It becomes dangerous when a low number is treated as permission to
stop reading.

\section{Pattern lists begin the diagnosis}
\label{sec:oss-humanizer}

Consider a deliberately small example:

\begin{repairpair}
\before{``The package contains a comprehensive suite of 35 robust patterns
that collectively underscore the importance of human-centred writing.''}
\after{``The package lists 35 patterns found in AI-assisted prose.''}
\end{repairpair}

The revision is easier to read because the package now performs a concrete
action: it lists patterns. The number remains. Three unsupported evaluations,
``comprehensive,'' ``robust,'' and ``underscore the importance,'' disappear.
An editor can explain every change without claiming to know who wrote the
sentence.

Now imagine that the next sentence names all 35 patterns, divides them into
six classes, gives examples from four classes, states two exceptions, and
reports the package's licence and release status. It might contain none of the
phrases in the first sentence. It would still be hard to read. The writer has
opened several questions before closing the first one.

That distinction matters because blader/humanizer, the strongest general
pattern list we found, is deliberately a catalogue. It turns Wikipedia's
``Signs of AI writing'' into 35 editing prompts. Some prompts concern inflated
importance or vague sources. Others concern predictable sentence shapes,
headings, filler, and remnants of a chat exchange. Each prompt comes with an
example. The list is broad enough to give an editor a useful first look, and
the project already uses a pinned copy.

The list works best when it changes a suspicion into a question. If a passage
says that a result ``marks a pivotal moment,'' the editor can ask what changed
and by how much. If three consecutive paragraphs end with a broad outlook, the
editor can ask whether any of those conclusions adds evidence. If a sentence
uses ``experts say,'' the editor can ask which experts and where the claim can
be checked. The pattern points to a place; the author supplies the answer.

\hvFigurePatternQuestion

This is a less glamorous role than automatic humanization, but it is much more
valuable. The tool does not need to infer a writer's identity or invent a
better sentence. It needs to locate a plausible problem, name the reason, and
leave enough context for a person to decide.

The word ``significant'' shows why context comes next. In ``the coefficient is
significant at the one per cent level,'' the word has a stated statistical
meaning. In ``the project makes a significant contribution,'' it is an
evaluation waiting for evidence. A flat search returns both sentences. A good
finding shows the surrounding clause and asks only about the second use unless
the statistical statement itself lacks a reported test.

Two parts of humanizer make that use possible. First, it tells the editor what
not to flag. Formal vocabulary is not evidence of machine authorship. A
quotation belongs to the person being quoted. A genuine limitation stays when
the evidence requires it. An em dash, by itself, proves nothing. These
exceptions keep an ordinary feature of scholarly prose from becoming an
offence merely because it also appears in generated text.

\Needspace{11cm}
Second, the skill forbids a rewrite from changing a number, date, quotation,
citation, or substantive claim. Suppose a cited study estimates a lower bound
of 13.5 per cent. Changing the sentence to ``about 14 per cent'' may improve
the cadence, but it changes the reported result. A prose tool has no authority
to do that. It should either preserve 13.5 per cent or ask the author to review
the source.

The distinction is easy to miss because the altered sentence may sound more
natural. Fact preservation therefore has to be checked after the prose edit,
not entrusted to the editor's memory. The comparison should display 13.5 and
14 side by side, retain the citation attached to the number, and require an
author to accept or reject the change. Fluency is not evidence that rounding
was licensed.

Humanizer also contains a rule that illustrates the need for local control. It
discourages em and en dashes unless the writer's sample uses them. Applied
literally to LaTeX, that instruction would damage correct typography in a date
range such as 1990--2020 and in a name such as Nash--Cournot. Humanvoice keeps
the diagnostic but overrides the blanket instruction. The detailed record in
Appendix~\ref{app:oss-humanizer-record} preserves the version and implementation
facts; the reader-facing lesson is simply that a useful list must be allowed
to be wrong in context.

\section{Pace is not the same as brevity}
\label{sec:oss-pace}

A page can be spacious and still move too fast. Suppose a package review tells
the reader, in one stretch, that the project contains 35 patterns in six
groups; names examples from each group; explains its false-positive policy and
its fact-preservation rule; identifies a dangerous typography rule; describes
the local override; and reports that the vendored copy matches upstream. The
sentences may be grammatical. The paragraph may contain fewer than 350 words.
It is nevertheless doing the work of several pages.

The reader has to keep too many questions open. What does the package actually
catch? Why are the six groups useful? Which rule is dangerous? Why does a
typographic exception affect adoption? What has already been installed, and
what remains to be built? The answers arrive, but they arrive while the reader
is still sorting the questions.

This is the kind of failure a word count misses. It is also easy to worsen by
shortening. Removing connective sentences leaves the same ideas packed more
tightly. Breaking the paragraph into bullets changes its shape but not its
pace. Adding a diagram can make matters worse if the diagram introduces a new
taxonomy beside an already crowded explanation.

The CardNPV case in Section~\ref{sec:cardnpv} uses the more reliable sequence.
It begins with one finite account and one decision. The reader sees the entries,
follows the arithmetic, and learns why a particular timing distinction matters.
Only then does the text give the distinction a general name and carry it to a
larger model. The example does not decorate an abstraction that has already
been stated. It creates the need for the abstraction.

The same sequence works for software. Begin with one sentence that a tool can
plausibly improve. Show the pattern it notices. Then show a sentence on which
the same rule would be wrong. Once the reader understands both cases, explain
how the product will use the rule. Repository status and licence can wait until
the reader knows why the package matters; the complete facts can sit in a
reference record.

\hvFigureConceptPace

Slower exposition need not be repetitive. Each return should add something.
The first example establishes the symptom. The counterexample establishes the
boundary. A worked product finding shows the action. The implementation
decision then follows from objects the reader already knows. Repeating the name
of a category without changing the reader's understanding would merely add
length.

Pace also changes at natural stopping points. An equation should usually be
followed by an interpretation before another notation system appears. A table
should be introduced by the question it answers and followed by the one or two
comparisons that matter. A new package should not arrive while the implication
of the previous package remains unstated. These are ordinary teaching habits,
but a catalogue-driven survey tends to erase them because every record competes
for the same small template.

Humanvoice should therefore diagnose crowding separately from sentence style.
A first implementation can nominate passages with several new acronyms,
package names, parenthetical qualifications, enumerations, or cross-references
in quick succession. It can also notice when a short section introduces many
terms that do not recur. Those are reasons to inspect a page, not evidence that
the page is bad. A compact theorem statement may be exactly right for its
reader; a compact opening to a product proposal usually is not.

The decisive observation still comes from a person. Give a reader the page and
ask for the relation among its principal ideas. If the reader can repeat the
nouns but cannot explain why one leads to the next, the problem is not
vocabulary. The page has named a structure without teaching it. Revision should
then change the order, examples, and connections before it changes the tone.

\section{Context decides whether a rule applies}
\label{sec:oss-context}

Take the sentence ``The estimator is asymptotically normal.'' A general style
guide that removes every adverb would delete the word carrying the mathematical
content. The shorter sentence would be false. A ban on passive voice can fail
in the same way. ``The observations were winsorized at the first and
ninety-ninth percentiles'' puts the treatment first because the identity of the
research assistant is irrelevant. Recasting it merely to name an actor can make
the method less clear.

These are not edge cases in technical writing. They are the reason a package
written for essays, blog posts, or sales copy cannot be installed as a
scholarly rulebook. hardikpandya/stop-slop, for example, notices several habits
that our own drafts exhibit: distant narration, staged negative openings, and
sentences in which an abstraction appears to act. Those observations are
useful. Its flat prohibitions on adverbs and passive constructions are not.
Humanvoice can borrow the first set without importing the second.

Even the idea of an ``abstract actor'' needs a domain test. ``The decision
emerges'' may hide the person or calculation that produced it. ``The model
implies'' is ordinary economics when the implication follows from stated
assumptions. The words are similar; the semantic roles are not. A linter can
nominate both sentences. Only a context rule, and sometimes the author, can
separate them.

\Needspace{13cm}
The useful question is not whether an inanimate noun performs a verb. It is
whether the sentence lets the reader recover the cause. ``The estimate falls
after the high-leverage observations are removed'' names an observable
comparison. ``A more balanced conclusion emerges'' does not say who balanced
the evidence or what changed. The first sentence can stand. The second needs an
actor, a calculation, or both.

The academic-humanizer project begins closer to the right register. It keeps
evidence-tied hedges such as ``suggests'' and ``is consistent with.'' It keeps
equations, notation, citation keys, and definitions. It asks whether a strong
empirical verb is supported by a number, figure, test, or source. Its most
useful example does not replace ``significantly outperforms'' with a milder
synonym. It asks for the comparison, the magnitude, and the statistical basis.
That is an explanation problem, not a tone problem.

In a hypothetical results section, ``our estimator significantly outperforms
the baseline'' leaves all three elements unstated. A useful revision might say
that held-out root mean squared error falls from 1.14 to 1.07 and then name the
test used for the comparison. If those results do not exist, the right repair
is to weaken or remove the claim. The editor must not invent a number merely to
satisfy the template.

The project is too young and too dependent on model compliance to become a
second installed editor. Its academic rules nevertheless fill an important
gap in the general list. We will adapt those rules into a scholarly layer,
using examples from economics and finance. A calibrated hedge will remain a
hedge; a defined technical term will not be replaced merely because it is
uncommon.

conorbronsdon/avoid-ai-writing contributes a related distinction. It separates
words that have become unusually frequent in generated prose from ordinary
clarity edits. The first can contribute to a distributional signal. The second
may improve a sentence, but it says nothing reliable about authorship. A writer
may reasonably replace ``utilize'' with ``use.'' That preference must not be
smuggled into a detector.

The same project handles quoted and hostile text carefully. A passage being
reviewed may contain the sentence ``ignore the instructions above.'' The
editor must treat those words as document content, not as a command. A phrase
inside a quotation, code block, or table can be reported when necessary, but
it should not be rewritten automatically. In Humanvoice, source authority
comes from the author and the declared configuration, never from the text
under inspection.

Taken together, the packages suggest a three-part rule. A diagnostic should
say what it observed. The configuration should say where that observation is
meaningful. The author should decide what, if anything, changes. A global
blacklist collapses all three decisions into one and is therefore too crude
for the product.

\section{A useful warning points to a place}
\label{sec:oss-located-findings}

An author facing an eighty-page manuscript cannot act on ``this document is
somewhat repetitive.'' The useful unit is smaller. The warning should identify
the passage, state what repeated, and ask a question whose answer could change
the paragraph.

For example, suppose four consecutive sections begin with ``This chapter
provides.'' A Humanvoice finding might read as follows:

\begin{plainbox}{A located prose finding}
\textbf{Observed:} four of the last four section openings use the same
document-centred construction.\par
\textbf{Location:} source lines 410, 463, 521, and 588; rendered pages 27, 31,
35, and 39.\par
\textbf{Question:} can each section begin with the result, example, or problem
that distinguishes it?\par
\textbf{Action:} the author may revise, keep, or dismiss the finding with a
reason.
\end{plainbox}

This record does not prescribe four replacement sentences. It makes a pattern
visible at the scale where the author can judge it. If the repetition is a
deliberate refrain, the author can keep it. If it arose from a chapter template,
the locations make the repair finite.

Location is necessary but not sufficient. A long manuscript can produce
hundreds of accurate low-value warnings. Humanvoice should present first the
findings that could change a reader's interpretation: a missing qualification,
an unsupported comparison, a broken referent, or a repeated section structure
that hides the argument. Spelling and punctuation remain available, but they
should not fill the screen while a substantive problem waits below them.

This ordering is not a universal severity formula. A misspelled unit in a risk
limit can matter more than an awkward paragraph. The finding record therefore
retains its rule, location, and affected object. The reading brief supplies the
priority: what must this reader understand or decide, and which warning bears
on that task?

Vale is the strongest available host for findings of this kind. Given an
explicit rule, it returns a file, line, severity, and message. Its configuration
can turn a rule off for LaTeX while retaining it for a press release, or reduce
an uncertain pattern from an error to advice. That is how genre and policy
become executable without pretending that every rule is universal.

The accompanying vale-ai-tells collection supplies useful raw material. It
contains ordinary vocabulary rules, patterns for defensive sentence shapes,
and measurements of such features as sentence-opener repetition and paragraph
variation. Several stock rules are wrong for scholarly LaTeX. One treats every
en dash as an error; another flags the double hyphen that LaTeX uses to produce
one. A formal-register rule can even object to ``minimize,'' though an estimator
may literally minimize an objective. The remedy is a small, reviewed
configuration, not a mass acceptance of the package defaults.

sloptrim complements Vale. Vale is good at saying which explicit rule fired on
which line. sloptrim gives a quick view of recurring vocabulary and rhythm
across the document, with sample spans that can be inspected. It also publishes
its own limits and does not claim to identify authorship. That candour matters:
the score in Table~\ref{tab:calibration} is useful precisely because we do not
ask it to certify the page.

The two tools should not be merged into one opaque score. If the document-level
scan reports monotonous openings, the author should be able to open the spans
that produced the count. If a Vale rule fires on one of those spans, the record
should show the rule separately. Agreement makes the passage easier to find;
disagreement reveals that the tools are measuring different things.

LaTeX exposes one more wrinkle. A program reading raw source may mistake a
sequence such as \verb|\item \textbf| for a repeated sentence opening. A
program reading only extracted prose may lose the source line that an author
needs to fix. The first implementation should therefore retain both views. Raw
source preserves locations and commands. Extracted prose gives the diagnostic
a cleaner linguistic sample. When they disagree, the finding should show the
disagreement rather than silently choosing one.

\hvFigureDualTextViews

This division yields a modest first stack. Vale supplies located, configurable
rules. sloptrim supplies a reproducible document-level scan. The installed
humanizer and the adapted academic layer supply questions and examples. None
of them rewrites by default, and none of them closes the reader test.

\section{LaTeX changes the engineering problem}
\label{sec:oss-latex-problem}

A prose linter sees words. A technical revision also contains equations,
numbers, citations, labels, tables, and commands that acquire meaning only
when LaTeX builds the document. Consider these two expressions:

\begin{align*}
\widehat\beta_A &= (X^{\mathsf T}X)^{-1}X^{\mathsf T}y,\\
\widehat\beta_B &= (XX^{\mathsf T})^{-1}X^{\mathsf T}y.
\end{align*}

One transposition has moved. A style checker may report no change at all. A
character diff will show the change but cannot say whether it is admissible.
The author needs the equation highlighted as a protected difference, with its
label and rendered location, before deciding whether version B is an error or
an intentional change of estimator.

A protected comparison handles this in stages. It first fixes the identity of
the two source versions. It then records the complete equation bodies and their
labels, rather than comparing only the prose around them. After both versions
compile, it connects each source expression to the page on which the reader
sees it. Only then does it ask the author about the moved transposition. The
software establishes where the change occurred; the author establishes what
the change means.

The same problem appears in less conspicuous forms. A citation key can still
compile after it has been attached to the wrong claim. A custom command can
expand to a number in one table and a label in another. A revised paragraph can
leave an equation untouched while deleting the assumption that made it valid.
No prose package reviewed here represents all of those relations.

\Needspace{11cm}
Consider the citation case. A revision can replace one source with another
whose title and date look plausible. LaTeX will compile as long as the new key
exists. A bibliography resolver can establish the new source's identity. It
still cannot establish that the source supports the sentence. Humanvoice must
show the changed identity beside the claim and leave support checking to a
person or a separately evidenced process.

TeXtidote and LTeX+ are useful because they understand enough LaTeX to skip
mathematics and return grammar findings to source lines. textlint offers a more
programmable rule system and a LaTeX parser plug-in, but would require us to
write the relevant rules. These are candidates for line-accurate language
feedback. They are not substitutes for a protected comparison.

The deeper source question belongs to the next chapter. Tools such as
pylatexenc, Pandoc, plasTeX, unified-latex, TexSoup, tree-sitter-latex, and
LaTeXML offer different compromises among tolerance, macro expansion, source
locations, and document structure. Their names matter to an implementer. To a
reader of the proposal, the decision is simpler: benchmark several on the same
small set of hostile LaTeX examples, preserve raw spans, and abstain whenever a
parser does not understand a construct. Chapter~\ref{ch:software-dossiers}
develops that comparison at an engineering pace.

This is the point at which a writing assistant becomes a document system. The
software must connect a warning in prose to the exact source version, the
rendered page, and any mathematical or evidentiary object that could be changed
by the repair. Existing packages supply parts of that path. Humanvoice must
build and test the connection.

\section{Detection is not the product}
\label{sec:oss-detection-boundary}

It is tempting to turn the same tools into an answer to a more dramatic
question: was this document written by a model? That question is poorly matched
to the product and poorly supported at the level of an individual technical
manuscript.

The literature in Chapter~\ref{sec:literature} shows why. Detectors are
sensitive to model family, domain, editing, language background, and the
polished regularity of formal prose. A false accusation can attach to a person;
a missed label can be defeated by ordinary editing. Neither outcome tells an
executive whether the document explains its mechanism.

fast-detect-gpt remains useful as research code for a distributional measure.
If the team later studies a large collection of its own drafts, it may compare
the curvature distribution across periods or workflows. The observation must
remain aggregate, with no individual label and no claim of quality. The
experiment is optional; the writing workflow does not depend on it.

slop-forensics answers a more practical corpus question. It compares word and
phrase frequencies in model output with a human reference collection. Pointed
at the team's own drafts and a pre-2022 economics corpus, it could reveal that
a new model has begun to overproduce a particular transition or sentence
shape. The result would update the diagnostic list. It would not trigger an
automatic rewrite, because an overrepresented expression can still be the
right expression in a particular sentence.

The boundary also excludes tools organized around detector evasion. One
reviewed project describes its checks in terms of watermark stripping and
attribution evasion, and several of those checks pass without measuring the
property named. That goal conflicts with Humanvoice. The product is intended
to improve disclosed, accountable technical work. It has no need to make
assistance harder to identify.

Corpus monitoring may eventually help prevent all the team's documents from
acquiring the same voice. It is a population diagnostic, much like checking a
portfolio for concentration. The decision on an individual paragraph remains
local: does this wording convey the right fact, in the right register, to this
reader?

\section{The software decision}
\label{sec:oss-decision}

The survey supports four different actions, and keeping them separate prevents
an interesting repository from becoming an unnecessary dependency.

We will keep the installed humanizer as a general source of questions. Its
false-positive guidance and fact-preservation rule make it safer than a simple
ban list. We will adapt the strongest academic rules into the same policy
layer, so that evidence, notation, and calibrated hedging receive explicit
protection.

We will pilot Vale and sloptrim as complementary diagnostics. Vale must earn
its place by returning useful, line-addressable findings under a restrained
LaTeX configuration. sloptrim must earn its place by remaining stable on a
held-out set of accepted and rejected passages. A favorable score from either
tool will never count as reader acceptance.

We will test, rather than prejudge, the source components. LaTeX parsers,
language servers, and visual comparison tools will face the same fixtures:
custom commands, equations, tables, citation changes, malformed environments,
and intentional repetition. A parser that abstains visibly is preferable to
one that loses a protected object while producing a clean-looking tree.

We will leave out packages whose genre assumptions conflict with scholarly
writing, whose checks are weaker than available alternatives, or whose purpose
is detector evasion. Their useful observations remain in the pattern catalogue
or the package records. Their code does not enter the first build.

The field therefore supplies much of the machinery, but not the product. It
can parse parts of a source file, locate some prose patterns, compare visible
changes, and run a model locally. It does not join those observations into a
protected revision that ends with a named reader's decision. Chapter
\ref{ch:software-dossiers} now turns from lessons to components and asks which
ones can survive a held-out engineering test.
